\documentstyle[%


righttag%


,11pt%


,amssymb%


,russian%               remove in Eng.


,izvru%                 Set izven% in Eng.


]{amsart}





% 





\newcommand{\re}{\operatorname{Re}}


\newcommand{\im}{\operatorname{Im}}


\newcommand{\tts}[1]{}





%\hoffset=-15mm%                  For Eng.


\setcounter{page}{57}


\god{2000}


\nomer{10~(461)} %                        Remove in Eng.


%\nomer{Vol.\,44, No.\,10}%               Set in Eng.


%\str{\pageref{begin}--\pageref{end}}%  Set in Eng.


\udk{517.958\,:\,539.3}





\author{\it Н.Б.\,Плещинский, А.А.\,Гусенкова}


% NB:   ^^ remove in Eng.


\title{Комплексные потенциалы


с логарифмическими особенностями в ядрах


для упругих тел с дефектом вдоль гладкой дуги}





\date{}


%\pagestyle{myheadings}%         Set in Eng.


\pagestyle{plain} %              Remove in Eng.


\begin{document}


%\thispagestyle{plain}%          Set in Eng.


\maketitle


\label{begin}








В плоской теории упругости при отсутствии массовых (объемных) сил


напряжения и перемещения могут быть выражены через {\it две аналитические


функции} по формулам Колосова--Мусхелишвили (\cite{Mush.MTU}, \S\,32).


Интегральное уравнение, эквивалентное исходной граничной задаче,


может быть получено, если представить эти функции в виде интегралов


по границе рассматриваемой области и подставить их в граничные условия.





В данной работе построены {\it интегральные представления с


логарифмическими особенностями в ядрах} для функций, аналитических


в различных областях комплексной плоскости с разрезами вдоль гладкой дуги.


Рассмотрены полная плоскость, полуплоскость, круг и области, конформно


эквивалентные кругу (полуплоскости).


Получены {\it сингулярные интегральные уравнения} {\it с логарифмическими


особенностями в ядрах}, {\it эквивалентные первой и второй основным


граничным задачам} плоской теории упругости для тел с дефектами


(с разрезами или тонкими включениями).


Уточнены и дополнены результаты работы \cite{PNB.Pre}.





Пусть $ab$ --- гладкая разомкнутая дуга на комплексной плоскости.


Для функций, определенных в плоскости с разрезом по $ab$,


будем обозначать знаками $+$ и $-$ их предельные значения


на левом и правом (по отношению к направлению движения от $a$ к $b$)


берегах разреза.





Следуя принципам математической теории трещин, построенной в работах


В.В.Па\-на\-сю\-ка, М.П.Сав\-ру\-ка и др. (напр., \cite{Svr.DZUTT}),


назовем {\it комплексными потенциалами для упругого тела с дефектом


вдоль дуги} $ab$ функции $\varphi(\cdot)$ и $\psi(\cdot)$, аналитические в


соответствующей области комплексной плоскости с разрезом,


предельные значения которых на $ab$ удовлетворяют условиям


\begin{gather}


[\varphi(t)+t\overline{\varphi'(t)}+\overline{\psi(t)}]^+


- [\varphi(t)+t\overline{\varphi'(t)}+\overline{\psi(t)}]^- =


(k+1)\int_{at}\alpha(\tau)d\tau,


\quad t\in ab, \\


%


[k\varphi(t)-t\overline{\varphi'(t)}-\overline{\psi(t)}]^+


- [k\varphi(t)-t\overline{\varphi'(t)}-\overline{\psi(t)}]^- =


(k+1)\int_{at}\beta(\tau)d\tau,


\quad t\in ab,


\end{gather}


где $k>0$ --- некоторая вещественная постоянная ($k=3-4\mu$ в случае


плоской деформации или $k=(3-\mu)/(1+\mu)$ в случае обобщенного плоского


напряженного состояния, $\mu$ --- коэффициент Пуассона)


и $\alpha(\cdot),\;\beta(\cdot)$ ---


некоторые определенные на $ab$ функции, причем


\begin{equation}


\int_{ab}\alpha(\tau)d\tau=0, \quad \int_{ab}\beta(\tau)d\tau=0.


\end{equation}


Будем считать, что функции $\alpha(\cdot),\;\beta(\cdot)$ удовлетворяют


условию Г\"ельдера на любой части дуги $ab$, не содержащей концов,


а в точках $a$ и $b$ могут иметь особенности интегрируемого порядка


(принадлежат классу $H^*$ по Н.И.\,Мусхелишвили).





Условия (1) и (2) задают на $ab$ скачки напряжений и перемещений,


множитель $k+1$ в правой части используется для упрощения некоторых


формул.





\section{Упругая плоскость c дефектом вдоль гладкой дуги}





Пусть рассматриваемые в разрезанных по $ab$ областях аналитические


функции непрерывно продолжимы на все точки дуги, кроме, может быть,


ее концов.





\begin{lem}


Если аналитические в разрезанной по дуге $ab$ комплексной плоскости


функции $\varphi(\cdot)$, $\psi(\cdot)$ ограничены на бесконечности и


удовлетворяют


условиям $(1)$--$(3)$, то


\begin{gather}


\varphi(z) = \frac{1}{2\pi i}\int_{ab} [\alpha(\tau)+\beta(\tau)]


\int_{\tau b}\frac{d\xi}{\xi-z}d\tau + P, \\


%


\psi(z) = - \frac{1}{2\pi i}\int_{ab} [\alpha(\tau)+\beta(\tau)]


\frac{\overline{\tau}d\tau}{\tau-z} +


\frac{1}{2\pi i}\int_{ab} [k\overline{\alpha(\tau)}


-\overline{\beta(\tau)}]


\int_{\tau b}\frac{d\xi}{\xi-z}\,d\overline{\tau} + Q,


\end{gather}


где $P$ и $Q$ --- произвольные комплексные постоянные.


\end{lem}





\begin{pf}


Сложив (1) и (2), получим


$$


\varphi^+(t)-\varphi^-(t)=\int_{at}[\alpha(\tau)+\beta(\tau)]d\tau,


\quad t\in ab.


$$


Решение этой задачи о скачке дает формула (4). Тогда, учитывая


равенства (3), найдем


$$


\varphi'(z) = \frac{1}{2\pi i}\int_{ab} [\alpha(\tau)+\beta(\tau)]


\frac{d\tau}{\tau-z},\quad


\varphi'{}^+(t)-\varphi'{}^-(t) = \alpha(t)+\beta(t),\quad t\in ab.


$$


Умножив равенство (1) на $k$ и вычитая из него (2), получим


$$


\overline{\psi^+(t)} - \overline{\psi^-(t)} =


\int_{at} [k\alpha(\tau)-\beta(\tau)]\;d\tau


- t\,[\overline{\varphi'{}^+(t)} - \overline{\varphi'{}^-(t)}],


\quad t\in a b,


$$


а отсюда следует представление (5).


\end{pf}





Если задать значения функций $\varphi(\cdot)$ и $\psi(\cdot)$


в некоторой точке $z_0$, то постоянные $P$ и $Q$ однозначно определяются.


Например, если $z_0=\infty$, то $P=\varphi(\infty)$, $Q=\psi(\infty)$.





\begin{cor}


Если функции $\varphi(\cdot),\;\psi(\cdot)$ удовлетворяют условиям леммы 1,


то


\begin{gather}


[\varphi(t)+t\overline{\varphi'(t)}+\overline{\psi(t)}]^+


+ [\varphi(t)+t\overline{\varphi'(t)}+\overline{\psi(t)}]^- = \notag \\


%


= \frac{1}{\pi i}\int_{ab} [\alpha(\tau)+\beta(\tau)]


\int_{\tau b}\frac{d\xi}{\xi-t}\,d\tau


- \frac{1}{\pi i}\int_{ab} [k\alpha(\tau)-\beta(\tau)]


\int_{\tau b}\frac{d\overline{\xi}}{\overline{\xi}-\overline{t}}\,


d\tau + \notag \\


%


+ \frac{1}{\pi i}\int_{ab} [\overline{\alpha(\tau)}


+\overline{\beta(\tau)}]


\frac{\tau-t}{\overline{\tau}-\overline{t}}\;d\overline{\tau} + 2P


+ 2\overline{Q}, \\


%


[k\varphi(t)-t\overline{\varphi'(t)}-\overline{\psi(t)}]^+


+ [k\varphi(t)-t\overline{\varphi'(t)}-\overline{\psi(t)}]^- = \notag \\


%


= \frac{k}{\pi i}\int_{ab} [\alpha(\tau)+\beta(\tau)]


\int_{\tau b}\frac{d\xi}{\xi-t}\,d\tau


+ \frac{1}{\pi i}\int_{ab} [k\alpha(\tau)-\beta(\tau)]


\int_{\tau b}\frac{d\overline{\xi}}{\overline{\xi}-\overline{t}}\,


d\tau - \notag \\


%


- \frac{1}{\pi i}\int_{ab} [\overline{\alpha(\tau)}


+\overline{\beta(\tau)}]


\frac{\tau-t}{\overline{\tau}-\overline{t}}\,d\overline{\tau} + 2kP


- 2\overline{Q}.


\end{gather}


\end{cor}





Рассмотрим первую и вторую основные граничные задачи для упругой


плоскости с дефектом вдоль дуги $ab$, когда на берегах разреза заданы


напряжения (несамоуравновешенные усилия)


$$


[\varphi(t)+t\overline{\varphi'(t)}+\overline{\psi(t)}]^{\pm} =


\int_{at}f^{\pm}(\tau)d\tau + A,\quad t\in ab,


$$


или перемещения (в интегральной форме)


$$


[k\varphi(t)-t\overline{\varphi'(t)}-\overline{\psi(t)}]^{\pm} =


\int_{at}g^{\pm}(\tau)d\tau + B, \quad t\in ab,


$$


здесь постоянные $A$ и $B$ определяют


напряжение и перемещение в точке $a$ дуги.





С помощью формул (6) и (7) легко получить следующее утверждение.





\begin{thm}


Первая и вторая основные граничные задачи для упругой плоскости


с дефектом вдоль дуги $ab$ эквивалентны интегральным уравнениям


\begin{multline}


\frac{1}{\pi i}\int_{ab} [\alpha(\tau)+\beta(\tau)]


\int_{\tau b}\frac{d\xi}{\xi-t}\,d\tau


- \frac{1}{\pi i}\int_{ab} [k\alpha(\tau)-\beta(\tau)]


\int_{\tau b}\frac{d\overline{\xi}}{\overline{\xi}


-\overline{t}}d\tau + \\


%


+ \frac{1}{\pi i}\int_{ab} [\overline{\alpha(\tau)}


+\overline{\beta(\tau)}]


\frac{\tau-t}{\overline{\tau}-\overline{t}}\;d\overline{\tau}


= \int_{at} [f^+(\tau)+f^-(\tau)]d\tau + 2(A-P-\overline{Q}),


\quad t\in ab,


\end{multline}


где $\alpha(t)=f^+(t)-f^-(t)$, а $\beta(\cdot)$ -- искомая функция, и


\begin{multline}


\frac{k}{\pi i}\int_{ab} [\alpha(\tau)+\beta(\tau)]


\int_{\tau b}\frac{d\xi}{\xi-t}\,d\tau


+ \frac{1}{\pi i}\int_{ab} [k\alpha(\tau)-\beta(\tau)]


\int_{\tau b}\frac{d\overline{\xi}}{\overline{\xi}-\overline{t}}\,


d\tau - \\


%


- \frac{1}{\pi i}\int_{ab} [\overline{\alpha(\tau)}


+\overline{\beta(\tau)}]


\frac{\tau-t}{\overline{\tau}-\overline{t}}\,d\overline{\tau}


= \int_{at} [g^+(\tau)+g^-(\tau)]d\tau + 2(B-kP+\overline{Q}),


\quad t\in ab ,


\end{multline}


где $\beta(t)=g^+(t)-g^-(t)$, а $\alpha(\cdot)$ --- искомая функция.


\end{thm}





Интегральные уравнения (8) и (9) являются аналогами


уравнений (I.78) и (I.81) из \cite{Svr.DZUTT}. Отличие


состоит в том, что ядра уравнений (I.78) и (I.81) содержат особенности


1-го порядка, а ядра уравнений (8) и (9) --- особенности


логарифмического типа \cite{PNB.PTIU}, \cite{PNB.92}.


Интегральные уравнения работы \cite{Svr.DZUTT} могут быть получены


из уравнений (8) и (9) дифференцированием, что ухудшает их свойства.


Как показано в \cite{PNB.PTIU}, от уравнений (8), (9) легко


перейти к эквивалентным сингулярным интегральным уравнениям с ядром


Коши относительно первообразных искомых функций. При этом


$$


\int_{ab}\beta(\tau) \int_{\tau b} \frac{d\xi}{\xi-z}\,d\tau =


\int_{ab} \frac{\gamma(\xi)\,d\xi}{\xi-z},\quad


\gamma(\xi)=\int_{a\xi} \beta(\tau)d\tau .


$$


Такие сингулярные уравнения более устойчивы при численном решении


квадратурными методами типа метода механических квадратур.





Заметим, что первая основная граничная задача сводится к случаю,


когда $f^+(\cdot)=f^-(\cdot)$. Для доказательства используем прием, указанный


в (\cite{PP.IUTU}, \S\,23). Действительно, введем новую искомую функцию


$$


\psi_1(z)=\psi(z)-\overline{\psi_0(z)},\quad


\psi_0(z) = \frac{1}{2\pi i}\int_{ab}


[f^+(\tau)-f^-(\tau)]\ln\frac{b-z}{\tau-z}\,d\tau.


$$


Тогда


\begin{multline*}


[\varphi(t)+t\overline{\varphi'(t)}+\overline{\psi_1(t)}]^{\pm} =


[\varphi(t)+t\overline{\varphi'(t)}+\overline{\psi(t)}]^{\pm}


\mp\frac{1}{2}\int_{at}[f^+(\tau)-f^-(\tau)]d\tau - \\


%


- \frac{1}{2\pi i}\int_{ab}


[f^+(\tau)-f^-(\tau)]\ln\frac{b-t}{t-\tau}\,d\tau = \\


%


= \frac{1}{2\pi i}\int_{ab}


[f^+(\tau)-f^-(\tau)]\ln\frac{b-t}{t-\tau}\,d\tau


+ \frac{1}{2}\int_{at}[f^+(\tau)+f^-(\tau)]d\tau + A,


\end{multline*}


здесь выбрана ветвь многозначной логарифмической функции


$$


\ln\frac{b-t}{t-\tau} =\int_{\tau b}\frac{d\xi}{\xi-t}


\overset{\text{def}}{=} \frac{1}{2} \bigg(


\bigg[ \ln\frac{t-b}{t-\tau} \bigg]^+


+ \bigg[ \ln\frac{t-b}{t-\tau} \bigg]^- \bigg),\quad t\in (\tau b).


$$


Точно так же можно показать, что вторая граничная задача


сводится к случаю, когда $g^+(\cdot)=g^-(\cdot)$.





\section{Полуплоскость c дефектом вдоль дуги}





Пусть $D^+$, $D^-$ --- верхняя и нижняя полуплоскости комплексной


плоскости и $L$ --- вещественная ось.





\begin{lem}


Если аналитические в разрезанной по дуге $ab$ нижней полуплоскости $D^-$


функции $\varphi(\cdot)$, $\psi(\cdot)$


\begin{itemize}


\item[1)] непрерывно продолжимы на все точки вещественной оси, 


включая бесконечно удаленную точку, и


\begin{equation}


\varphi^-(t)+t\overline{\varphi'{}^-(t)}+


\overline{\psi^-(t)}=0, \quad t\in L,


\end{equation}


\item[2)] удовлетворяют условиям $(1)$--$(3)$,


\end{itemize}


то


\begin{align}


\varphi(z) &=


\frac{1}{2\pi i}\int_{ab} [\alpha(\tau)+\beta(\tau)]


\int_{\tau b}\frac{d\xi}{\xi-z}\,d\tau


+ \frac{1}{2\pi i}\int_{ab} [k\alpha(\tau)-\beta(\tau)]


\int_{\tau b}\frac{d\overline{\xi}}{\overline{\xi}-z}\,d\tau - \notag \\


%


&\qquad - \frac{1}{2\pi i}\int_{ab} [\overline{\alpha(\tau)}


+\overline{\beta(\tau)}]


\,\frac{\tau-\overline{\tau}}{\overline{\tau}-z}\,d\overline{\tau}+ R, \\


%


\psi(z) &=


- \frac{1}{2\pi i}\int_{ab} [\alpha(\tau)+\beta(\tau)]


\frac{\overline{\tau}}{\tau-z}\,d\tau


- \frac{1}{2\pi i}\int_{ab} [k\alpha(\tau)-\beta(\tau)]


\frac{z}{\overline{\tau}-z}\,d\tau + \notag \\


%


&\qquad + \frac{1}{2\pi i}\int_{ab} [\overline{\alpha(\tau)}


+\overline{\beta(\tau)}]


\bigg[\,


\int_{\tau b}\frac{d\overline{\xi}}{\overline{\xi}-z}


+ \frac{z(\tau-\overline{\tau})}{(\overline{\tau}-z)^2}


\bigg]\,d\overline{\tau} + \notag \\


%


&\qquad\qquad + \frac{1}{2\pi i}\int_{ab} [k\overline{\alpha(\tau)}


-\overline{\beta(\tau)}]


\int_{\tau b}\frac{d\xi}{\xi-z}\,d\overline{\tau}


- \overline{R},


\end{align}


где $R$ --- произвольная комплексная постоянная.


\end{lem}





\begin{pf}


Пусть $\varphi_1(\cdot)$, $\psi_1(\cdot)$ --- определенные


формулами (4), (5) потенциалы для разрезанной по дуге $ab$ плоскости.


Введем новые искомые функции


$\varphi_0(\cdot)=\varphi(\cdot)-\varphi_1(\cdot)$ и


$\psi_0(\cdot)=\psi(\cdot)-\psi_1(\cdot)$, для которых дуга $ab$ уже не


является особой


линией. Рассмотрим вспомогательную функцию, аналитическую в $D^+$ и $D^-$


и непрерывно продолжимую на $L$ сверху и снизу,


$$


F(z) = \big\{ \varphi_0(z),\ z\in D^-;


\ -z\overline{\varphi_0'(\overline{z})} -\overline{\psi_0(\overline{z})},


\ z\in D^+ \big\}.


$$


Отметим, что в нижней полуокрестности бесконечно удаленной точки


$\varphi_0(z)=c_0+c_1z^{-1}+c_2z^{-2}+\dotsb$. Поэтому существует конечный


предел произведения $z\overline{\varphi_0'(\overline{z})}$ при


$z\to\infty\in L$ по любому пути из $D^+$.


Разность предельных значений $F(\cdot)$ на оси в силу (10)


$$


F^+(t)-F^-(t)=f(t),\quad


f(t)=\varphi_1^-(t)+t\overline{\varphi_1'{}^-(t)}+\overline{\psi_1^-(t)},


\quad t\in L.


$$


Тогда, как решение задачи о скачке,


$$


F(z) = \frac{1}{2\pi i}\int_{-\infty}^{\infty}\frac{f(x)dx}{x-z}+C,


$$


где $C$ -- произвольная комплексная постоянная.


Подставив выражение функции


\begin{multline}


f(x) = \frac{1}{2\pi i}\int_{ab} [\alpha(\tau)+\beta(\tau)]


\int_{\tau b}\frac{d\xi}{\xi-x}\,d\tau


- \frac{1}{2\pi i}\int_{ab}[k\alpha(\tau)-\beta(\tau)]


\int_{\tau b}\frac{d\overline{\xi}}{\overline{\xi}-\overline{x}}\,


d\tau + \\


%


+ \frac{1}{2\pi i}\int_{ab} [\overline{\alpha(\tau)}


+\overline{\beta(\tau)}]


\,\frac{\tau-x}{\overline{\tau}-\overline{x}}\,d\overline{\tau} +P


+\overline{Q},\quad x\in L,


\end{multline}


в сингулярный интеграл, учтем, что $\overline{x}=x$, и изменим порядок


интегрирования. Методом вычетов легко вычислить интегралы


(удобно использовать готовую формулу из (\cite{Gah.KZ.77}, пример 26a


к главе 1))


\begin{align*}


\frac{1}{2\pi i}\int_{-\infty}^{\infty}


\frac{dx}{(\xi-x)(x-z)} &=


\bigg\{ \frac{1}{\xi-z},\ z\in D^+;


\ 0,\ z\in D^- \bigg\}, \\


%


\frac{1}{2\pi i}\int_{-\infty}^{\infty}


\frac{dx}{(\overline{\xi}-x)(x-z)} &=


\bigg\{0,\ z\in D^+;


\ - \frac{1}{\overline{\xi}-z},\ z\in D^- \bigg\}, \\


%


\frac{1}{2\pi i}\int_{-\infty}^{\infty}


\frac{dx}{x-z} &=


\bigg\{ \frac{1}{2},\ z\in D^+;


\ -\frac{1}{2},\ z\in D^- \bigg\}, \\


%


\frac{1}{2\pi i}\int_{-\infty}^{\infty}


\frac{\tau-x}{\overline{\tau}-x}\frac{dx}{x-z} &=


\bigg\{ \frac{1}{2},\ z\in D^+;


\ - \frac{1}{2} - \frac{\tau-\overline{\tau}}{\overline{\tau}-z},


\ z\in D^- \bigg\} .


\end{align*}


Тогда


\begin{align*}


F(z) &= \frac{1}{2\pi i}\int_{ab} [\alpha(\tau)+\beta(\tau)]


\int_{\tau b}\frac{d\xi}{\xi-z}\,d\tau + \frac{P+\overline{Q}}{2}+C,


\quad z\in D^+, \\


%


F(z) &= \frac{1}{2\pi i}\int_{ab} [k\alpha(\tau)-\beta(\tau)]


\int_{\tau b}\frac{d\overline{\xi}}{\overline{\xi}-z}\,d\tau - \\


%


&\qquad - \frac{1}{2\pi i}\int_{ab} [\overline{\alpha(\tau)}


+\overline{\beta(\tau)}]


\,\frac{\tau-\overline{\tau}}{\overline{\tau}-z}\,d\overline{\tau}


- \frac{P+\overline{Q}}{2}+C,


\quad z\in D^-,


\end{align*}


и


\begin{alignat*}{3}


\varphi(z)&=\varphi_0(z)+\varphi_1(z),&\quad


\varphi_0(z)&=F(z),&\quad z&\in D^-, \\


%


\psi(z)&=\psi_0(z)+\psi_1(z),&\quad


\psi_0(z)&=-zF'(z)-\overline{F(\overline{z})},&\quad z&\in D^-.


\end{alignat*}


Подставив сюда функцию $F(\cdot)$ при $z\in D^-$ и $z\in D^+$


и вместо $\varphi_1(z)$, $\psi_1(z)$ правые части формул (4) и (5),


получим (11) и (12), где $R=(P-\overline{Q})/2+C$.


Постоянные $P$ и $Q$ в выражениях потенциалов


$\varphi_1(\cdot)$ и $\psi_1(\cdot)$ можно было бы и не учитывать (считать их


равными нулю).





Постоянная $R$ определится, если задать значение


одной из функций $\varphi(\cdot)$ или $\psi(\cdot)$ в некоторой точке $z_0$


нижней полуплоскости или вещественной оси.


\end{pf}





\begin{cor}


Если аналитические в нижней полуплоскости функции


$\varphi(\cdot)$, $\psi(\cdot)$


непрерывно продолжимы на все точки вещественной оси, включая


бесконечно удаленную точку, и удовлетворяют условию (10),


то $\varphi(z)=C$, $\psi(z)=-\overline{C}$, где $C$ --- произвольная


комплексная постоянная. При этом


$\varphi(z)+\overline{\psi(z)}=0$\, $\forall z\in D^-\cup L$.


\end{cor}





\begin{thm}


Первая и вторая основные граничные задачи для упругой полуплоскости


с дефектом вдоль дуги $ab$ эквивалентны интегральным уравнениям


\begin{multline*}


\frac{1}{\pi i}\int_{ab} [\alpha(\tau)+\beta(\tau)]


\bigg[\,


\int_{\tau b}\frac{d\xi}{\xi-t} - \int_{\tau b}\frac{d\xi}{\xi


-\overline{t}}


+ \frac{(t-\overline{t})(\overline{\tau}-\tau)}{(\tau-\overline{t})^2}


\bigg]\,d\tau + \\


%


+ \frac{1}{\pi i}\int_{ab} [k\alpha(\tau)-\beta(\tau)]


\bigg[\,


\int_{\tau b}\frac{d\overline{\xi}}{\overline{\xi}-t}


- \int_{\tau b}\frac{d\overline{\xi}}{\overline{\xi}-\overline{t}}


\bigg]\,d\tau + \\


%


+ \frac{1}{\pi i}\int_{ab} [\overline{\alpha(\tau)}


+\overline{\beta(\tau)}]


\bigg[\,


\frac{\tau-t}{\overline{\tau}-\overline{t}} - \frac{\tau-t}{\overline{\tau}


-t}\bigg]\,d\overline{\tau}


- \frac{1}{\pi i}\int_{ab} [k\overline{\alpha(\tau)}


-\overline{\beta(\tau)}]


\,\frac{t-\overline{t}}{\tau-\overline{t}}\,d\overline{\tau} = \\


%


= \int_{at} [f^+(\tau)+f^-(\tau)]d\tau + 2A,\quad t\in ab ,


\end{multline*}


где $\alpha(t)=f^+(t)-f^-(t)$, а $\beta(\cdot)$ --- искомая функция, и


\begin{multline*}


\frac{1}{\pi i}\int_{ab} [\alpha(\tau)+\beta(\tau)]


\bigg[k\int_{\tau b}\frac{d\xi}{\xi-t}


+ \int_{\tau b}\frac{d\xi}{\xi-\overline{t}}


- \frac{(t-\overline{t})(\overline{\tau}-\tau)}{(\tau-\overline{t})^2}


\bigg]\,d\tau + \\


%


+ \frac{1}{\pi i}\int_{ab} [k\alpha(\tau)-\beta(\tau)]


\bigg[k\int_{\tau b}\frac{d\overline{\xi}}{\overline{\xi}-t}


+ \int_{\tau b}\frac{d\overline{\xi}}{\overline{\xi}-\overline{t}}


\bigg]\,d\tau - \\


%


- \frac{1}{\pi i}\int_{ab} [\overline{\alpha(\tau)}


+\overline{\beta(\tau)}]


\bigg[\frac{\tau-t}{\overline{\tau}-\overline{t}}


+ k \frac{\tau-t}{\overline{\tau}-t}


\bigg]\,d\overline{\tau}


+ \frac{1}{\pi i}\int_{ab} [k\overline{\alpha(\tau)}


-\overline{\beta(\tau)}]


\,\frac{t-\overline{t}}{\tau-\overline{t}}\,d\overline{\tau} = \\


%


= \int_{at} [g^+(\tau)+g^-(\tau)]d\tau + 2B - 2(k+1)R,\quad t\in ab ,


\end{multline*}


где $\beta(t)=g^+(t)-g^-(t)$, а $\alpha(\cdot)$ --- искомая функция.


\end{thm}





Для доказательства достаточно вычислить на дуге $ab$ выражения


вида (6) и (7).





\section{Упругий круг c дефектом вдоль гладкой дуги}





Пусть теперь $D^+$ и $D^-$ --- круг $|z|<R$ и внешность круга


$|z|>R$, а $L$ --- окружность $|z|=R$.





\begin{lem}


Если аналитические в разрезанном по дуге $ab$ круге $D^+$


функции $\varphi(\cdot)$, $\psi(\cdot)$


\begin{itemize}


\item[1)] непрерывно продолжимы на все точки $L$ и


\begin{equation}


\varphi^+(t)+t\overline{\varphi'{}^+(t)}+


\overline{\psi^+(t)}=0, \quad t\in L,


\end{equation}


\item[2)] удовлетворяют условиям $(1)$--$(3)$,


\end{itemize}


то


\begin{align}


\varphi(z) &= \frac{1}{2\pi i} \int_{ab} [\alpha(\tau)+\beta(\tau)]


\int_{\tau b}\frac{d\xi}{\xi-z}\,d\tau + \notag \\


%


&\quad + \frac{1}{2\pi i} \int_{ab} [k\alpha(\tau)-\beta(\tau)]


\bigg[\,


\int_{\tau b}\frac{d\overline{\xi}}{\overline{\xi}-R^2/z}


- \frac{\overline{\tau}z}{2R^2}


\bigg]\,d\tau - \notag \\


%


&\quad\quad - \frac{1}{2\pi i} \int_{ab} [\overline{\alpha(\tau)}


+\overline{\beta(\tau)}]\bigg[


\frac{\tau-z}{\overline{\tau}-R^2/z} + \frac{\tau z}{2R^2}


\bigg]\,d\overline{\tau} + iSz-\overline{T} , \\


%


\psi(z) &= - \frac{1}{2\pi i} \int_{ab} [\alpha(\tau)+\beta(\tau)]


\frac{\overline{\tau}d\tau}{\tau-z}+ \notag \\


%


&\quad + \frac{1}{2\pi i} \int_{ab} [k\alpha(\tau)-\beta(\tau)]


\frac{(\overline{\tau}-\overline{b})


(\overline{\tau}\overline{b}z-\overline{\tau}R^2-\overline{b}R^2)}


{(\overline{\tau}z-R^2)(\overline{b}z-R^2)}d\tau +\notag \\


%


&\quad\quad + \frac{1}{2\pi i} \int_{ab} [\overline{\alpha(\tau)}


+\overline{\beta(\tau)}]\bigg[\,


\int_{\tau b}\frac{d\overline{\xi}}{\overline{\xi}-R^2/z}


+ \frac{(\tau\overline{\tau}-R^2)(\overline{\tau}z-2R^2)}


{(\overline{\tau}z-R^2)^2}


\bigg]\,d\overline{\tau} + \notag \\


%


&\quad\quad\quad + \frac{1}{2\pi i} \int_{ab} [k\overline{\alpha(\tau)}


-\overline{\beta(\tau)}]


\int_{\tau b}\frac{d\xi}{\xi-z}\,d\overline{\tau} + T ,


\end{align}


где $S$ и $T$ --- произвольные вещественная и комплексная постоянные.


\end{lem}





{\bf Доказательство.} Пусть, как и в случае полуплоскости,


$\varphi_1(\cdot)$, $\psi_1(\cdot)$


--- потенциалы для упругой плоскости с дефектом вдоль дуги $ab$, и


$\varphi_0(z)=\varphi(z)-\varphi_1(z)$ и $\psi_0(z)=\psi(z)-\psi_1(z)$ ---


новые искомые функции. Учитывая сказанное в конце доказательства леммы 2,


будем считать, что $P=Q=0$. Рассмотрим вспомогательную функцию


$$


F(z) = \big\{ \varphi_0(z),\ z\in D^+;


\ -z\overline{\varphi_0'(R^2/\overline{z})}


-\overline{\psi_0(R^2/\overline{z})},\ z\in D^- \big\},


$$


аналитическую в $D^+$ и $D^-$ всюду, кроме, может быть,


бесконечно удаленной точки, в окрестности которой поведение функции


$F(\cdot)$ определяет главная часть разложения в ряд Лорана


$-z\overline{\varphi'_0(0)}-\overline{\psi_0(0)}$. При этом


$$


F^+(t)-F^-(t)=-f(t),\quad


f(t)=\varphi_1^+(t)+t\overline{\varphi_1'{}^+(t)}+\overline{\psi_1^+(t)},


\quad t\in L.


$$


Тогда, как решение задачи о скачке,


$$


F(z)=F_0(z)-z\overline{\varphi'_0(0)}-\overline{\psi_0(0)},\quad


F_0(z) = -\,\frac{1}{2\pi i}\int_{L}\frac{f(x)dx}{x-z}.


$$


Так как $\varphi_0(z)=F(z),\, z\in D^+$, то при $z=0$


$$


\varphi_0(0)=F_0(0)-\overline{\psi_0(0)},\quad


\varphi'_0(0)=F'_0(0)-\overline{\varphi'_0(0)}.


$$


Поэтому $2\re \varphi'_0(0)=F'_0(0)$, а значение


$\im \varphi'_0(0)$


и одно из двух значений $\varphi_0(0)$ и $\psi_0(0)$ остаются произвольными.


Обозначим $S=\im \varphi'_0(0)$ и $T=\psi_0(0)$. Тогда


\begin{gather*}


F(z)=F_0(z)-\frac{z}{2}F'_0(0)+iSz-\overline{T},\quad z\in D^+\cup D^-, \\


%


\varphi_0(z)=F(z)=F_0(z)-\frac{z}{2}F'_0(0)+iSz-\overline{T},


\quad z\in D^+, \\


%


\psi_0(z)=-\frac{R^2}{z}\varphi_0'(z)


-\overline{F\bigg(\frac{R^2}{\overline{z}}\bigg)}


=-\frac{R^2}{z}\bigl[ F'_0(z)-F'_0(0)\bigl]


-\overline{F_0\bigg(\frac{R^2}{\overline{z}}\bigg)} +T,\quad z\in D^+


\end{gather*}


(отметим, что у функции $\psi_0(\cdot)$ полюса в точке $z=0$ нет).





Как и в случае полуплоскости, функция $f(\cdot)$ определена формулой


(13), но теперь $\overline{x}=R^2/x$. Вычислив интегралы по окружности $L$


\begin{align*}


\frac{1}{2\pi i}\int_{L}


\frac{dx}{(\xi-x)(x-z)} &=


\bigg\{ 0,\ z\in D^+; \ -\frac{1}{\xi-z},\ z\in D^- \bigg\}, \\


%


\frac{1}{2\pi i}\int_{L}


\frac{dx}{(\overline{\xi}-R^2/x)(x-z)} &=


\bigg\{\frac{z}{\overline{\xi}z-R^2},\ z\in D^+;


\ 0,\ z\in D^- \bigg\}, \\


%


\frac{1}{2\pi i}\int_{L}


\frac{\tau-x}{\overline{\tau}-R^2/x}\frac{dx}{x-z} &=


\bigg\{ \frac{z(\tau-z)}{\overline{\tau}z-R^2},\ z\in D^+;


\ 0,\ z\in D^- \bigg\} ,


\end{align*}


получим


\begin{align*}


F_0(z) &=


\frac{1}{2\pi i}\int_{ab} [k\alpha(\tau)-\beta(\tau)]


\int_{\tau b}\frac{zd\overline{\xi}}{\overline{\xi}z-R^2}\,d\tau - \\


%


&\qquad - \frac{1}{2\pi i}\int_{ab} [\overline{\alpha(\tau)}


+\overline{\beta(\tau)}]


 \,\frac{z(\tau-z)}{\overline{\tau}z-R^2}\,d\overline{\tau},


\quad z\in D^+, \\


%


F_0(z) &=


\frac{1}{2\pi i}\int_{ab} [\alpha(\tau)+\beta(\tau)]


\int_{\tau b}\frac{d\xi}{\xi-z}\,d\tau,


\quad z\in D^-. \quad\square


\end{align*}





\begin{cor}


Если аналитические в $D^+$ функции $\varphi(\cdot)$, $\psi(\cdot)$


непрерывно


продолжимы на $L$ и удовлетворяют условию (14), то


$\varphi(0)+\overline{\psi(0)}=0$, $\re\varphi'(0)=0$


и $\varphi(z)=iSz-\overline{T},\quad \psi(z)=T,$


где $S$ и $T$ --- произвольные вещественная и комплексная постоянные.


\end{cor}





Как и в случае полуплоскости, вычислив выражения


\begin{gather*}


[\varphi(t)+t\overline{\varphi'(t)}+\overline{\psi(t)}]^+


+ [\varphi(t)+t\overline{\varphi'(t)}+\overline{\psi(t)}]^-, \\


%


[k\varphi(t)-t\overline{\varphi'(t)}-\overline{\psi(t)}]^+


+ [k\varphi(t)-t\overline{\varphi'(t)}-\overline{\psi(t)}]^-,


\end{gather*}


можно доказать следующее утверждение.





\begin{thm}


Первая и вторая основные граничные задачи для упругого круга


с дефектом вдоль дуги $ab$ эквивалентны интегральным уравнениям


\begin{multline*}


\frac{1}{\pi i}\int_{ab} [\alpha(\tau)+\beta(\tau)]


\bigg[\,\int_{\tau b}\frac{d\xi}{\xi-t}


- \int_{\tau b}\frac{d\xi}{\xi-R^2/\overline{t}} - \\


%


- \frac{(\tau\overline{\tau}+t\overline{t})(\tau\overline{t}-2R^2)


+R^2(\overline{\tau}t-\tau\overline{t}+2R^2)}


{(\tau \overline{t}-R^2)^2}


+ \frac{\overline{\tau}t}{2R^2}\bigg]\,d\tau + \\


%


+ \frac{1}{\pi i}\int_{ab} [k\alpha(\tau)-\beta(\tau)]


\bigg[\,


\int_{\tau b}\frac{d\overline{\xi}}{\overline{\xi}-R^2/t}


- \int_{\tau b}\frac{d\overline{\xi}}{\overline{\xi}-\overline{t}}


- \frac{\overline{\tau}t}{2R^2}\bigg]\,d\tau + \\


%


+ \frac{1}{\pi i}\int_{ab} [\overline{\alpha(\tau)}


+\overline{\beta(\tau)}]


\bigg[\,


\frac{\tau-t}{\overline{\tau}-\overline{t}}


- \frac{\tau-t}{\overline{\tau}-R^2/t}


- \frac{\tau t}{2R^2}\bigg]\,d\overline{\tau} + \\


%


+ \frac{1}{\pi i}\int_{ab} [k\overline{\alpha(\tau)}


-\overline{\beta(\tau)}] \bigg[\,


\frac{tR^2}{\overline{t}(\tau\overline{t}-R^2)}


- \frac{(\tau-b)(\tau b\overline{t}-\tau R^2-bR^2)}


{(\tau\overline{t}-R^2)(b\overline{t}-R^2)}


+ \frac{\tau t}{2R^2}\bigg]\,d\overline{\tau} = \\


%


= \int_{at} [f^+(\tau)+f^-(\tau)]d\tau + 2A,\quad t\in ab ,


\end{multline*}


где $\alpha(t)=f^+(t)-f^-(t)$, а $\beta(\cdot)$ --- искомая функция, и


\begin{multline*}


\frac{1}{\pi i}\int_{ab} [\alpha(\tau)+\beta(\tau)]


\bigg[\, k\int_{\tau b}\frac{d\xi}{\xi-t}


+ \int_{\tau b}\frac{d\xi}{\xi-R^2/\overline{t}} + \\


%


+ \frac{(\tau\overline{\tau}+t\overline{t})(\tau\overline{t}-2R^2)


+R^2(\overline{\tau}t-\tau\overline{t}+2R^2)}


{(\tau \overline{t}-R^2)^2}


- \frac{\overline{\tau}t}{2R^2}\bigg]\,d\tau + \\


%


+ \frac{1}{\pi i}\int_{ab} [k\alpha(\tau)-\beta(\tau)] \bigg[\,


k\int_{\tau b}\frac{d\overline{\xi}}{\overline{\xi}-R^2/t}


+ \int_{\tau b}\frac{d\overline{\xi}}{\overline{\xi}-\overline{t}}


- k\frac{\overline{\tau}t}{2R^2} \bigg]\,d\tau - \\


%


- \frac{1}{\pi i}\int_{ab} [\overline{\alpha(\tau)}


+\overline{\beta(\tau)}] \bigg[\,


\frac{\tau-t}{\overline{\tau}-\overline{t}}


+ k\frac{\tau-t}{\overline{\tau}-R^2/t}


+ k\frac{\tau t}{2R^2} \bigg]\,d\overline{\tau} - \\


%


- \frac{1}{\pi i}\int_{ab} [k\overline{\alpha(\tau)}


-\overline{\beta(\tau)}] \bigg[\,


\frac{tR^2}{\overline{t}(\tau\overline{t}-R^2)}


- \frac{(\tau-b)(\tau b\overline{t}-\tau R^2-bR^2)}


{(\tau\overline{t}-R^2)(b\overline{t}-R^2)}


+ \frac{\tau t}{2R^2} \bigg]\,d\overline{\tau} = \\


%


= \int_{at} [g^+(\tau)+g^-(\tau)]d\tau + 2B + 2(k+1)(\overline{T}-iSt),


\quad t\in ab ,


\end{multline*}


где $\beta(t)=g^+(t)-g^-(t)$, а $\alpha(\cdot)$ --- искомая функция.


\end{thm}





\section{Области, конформно эквивалентные кругу (полуплоскости)}





Если упругая среда с дефектом вдоль разомкнутой дуги заполняет


область, конформно эквивалентную кругу (или полуплоскости), то комплексные


потенциалы также можно построить методом аналитического продолжения


через границу (методом сопряжения Н.И.\,Мусхелишвили).





Пусть области $D^+$ и $D^-$ расположены слева и справа от замкнутой


линии $L$ в плоскости комплексной переменной $z$, а области $\Delta^+$


и $\Delta^-$ --- слева и справа от линии $\Lambda$ в плоскости переменной


$\zeta$.


Рассмотрим параллельно два случая: 1)~когда $\Delta^+$ и $\Delta^-$ ---


единичный круг и внешность круга и 2)~когда $\Delta^+$ и $\Delta^-$ ---


верхняя и нижняя полуплоскости. 


Отметим, что оба случая эквивалентны.





\begin{lem}


Пусть рациональная функция $z=\omega(\zeta)$ взаимно однозначно отображает


$1)$~единичный круг $\Delta^+$ или $2)$~нижнюю 


полуплоскость $\Delta^-$ на область $D^+$.


Если предельные значения аналитических в разрезанной по дуге $ab$ области


$D^+$ функций $\varphi(\cdot)$, $\psi(\cdot)$ удовлетворяют условию


$$


\varphi^+(t)+t\overline{\varphi'{}^+(t)}+


\overline{\psi^+(t)}=0, \quad t\in L,


$$


и условиям $(1)$--$(3)$, то


\begin{equation}


\varphi(\omega(\zeta)) = \varphi_1(\omega(\zeta)) + \Phi_0(\zeta),\quad


\psi(\omega(\zeta)) = \psi_1(\omega(\zeta)) + \Psi_0(\zeta),


\end{equation}


где $\varphi_1(\cdot)$, $\psi_1(\cdot)$ -- комплексные потенциалы для


плоскости с дефектом вдоль $ab$,


\begin{align*}


&1)\quad \Phi_0(\zeta) = F(\zeta),\quad


\Psi_0(\zeta) = -\overline{F(1/\overline{\zeta})}


-\frac{\overline{\omega(1/\overline{\zeta})}}{\omega'(\zeta)}F'(\zeta),


\quad \zeta\in\Delta^+, \\


%


\intertext{или}


%


&2)\quad \Phi_0(\zeta) = F(\zeta),\quad


\Psi_0(\zeta) = -\overline{F(\overline{\zeta})}


-\frac{\overline{\omega(\overline{\zeta})}}{\omega'(\zeta)}F'(\zeta),


\quad \zeta\in\Delta^-,


\end{align*}


и $F(\cdot)$ --- решение задачи о скачке


\begin{equation}


\begin{array}{c}


F^+(\tau)-F^-(\tau) = \mp f(\omega(\tau)),\quad \tau\in\Lambda, \\[2mm]


%


f(t) = \varphi_1^+(t)+t\overline{\varphi_1'{}^+(t)}+


\overline{\psi_1^+(t)},\quad t\in L.


\end{array}


\end{equation}


\end{lem}





\begin{pf}


Пусть $\varphi_1(\cdot)$, $\psi_1(\cdot)$ --- комплексные


потенциалы (4), (5). Тогда предельные значения


на $L$ аналитических в $D^+$ функций


$\varphi_0(\cdot)=\varphi(\cdot)-\varphi_1(\cdot)$,


$\psi_0(\cdot)=\psi(\cdot)-\psi_1(\cdot)$ удовлетворяют условию


$$


\varphi_0^+(t)+t\overline{\varphi_0'{}^+(t)}+


\overline{\psi_0^+(t)} = -f(t),


\quad t\in L.


$$


Для новых искомых функций $\Phi_0(\zeta)=\varphi_0(\omega(\zeta))$,


$\Psi_0(\zeta)=\psi_0(\omega(\zeta))$, $\zeta\in\Delta^{\pm}$,


это условие примет вид


$$


\Phi_0^{\pm}(\tau)+ \frac{\omega(\tau)}{\overline{\omega'(\tau)}}


\overline{\Phi_0'{}^{\pm}(\tau)}


+\overline{\Psi_0^{\pm}(\tau)} = -f(\omega(\tau)),\quad \tau\in\Lambda .


$$


Верхний знак в формулах


соответствует случаю~1), а нижний --- случаю~2).


Рассмотрим вспомогательную функцию


\begin{align*}


&1)\quad F(\zeta) = \bigg\{ \Phi_0(\zeta),\ \zeta\in\Delta^+;


\ - \,\frac{\omega(\zeta)}{\overline{\omega'(1/\overline{\zeta})}}


\,\overline{\Phi_0'(1/\overline{\zeta})}


- \overline{\Psi_0(1/\overline{\zeta})},\ \zeta\in\Delta^- \bigg\} \\


%


\intertext{или}


%


&2) \quad F(\zeta) = \bigg\{ \Phi_0(\zeta),\ \zeta\in\Delta^-;


\ -\,\frac{\omega(\zeta)}{\overline{\omega'(\overline{\zeta})}} \,


\overline{\Phi_0'(\overline{\zeta})}


- \overline{\Psi_0(\overline{\zeta})},\ \zeta\in\Delta^+ \bigg\} .


\end{align*}


Можно показать, что $\omega'(\zeta)\ne 0$\, $\forall\zeta\in\Delta^{\pm}$


и $\omega'(\tau)\ne 0$\, $\forall\tau\in\Lambda$, если контур $\Lambda$


имеет непрерывно изменяющуюся кривизну (\cite{Mush.MTU}, \S\,47).


На единичной окружности $1/\overline{\tau}=\tau$, на вещественной оси


$\overline{\tau}=\tau$. Тогда предельные значения на $\Lambda$ функции


$F(\cdot)$ удовлетворяют условию (18).


\end{pf}





Решение задачи о скачке (18) должно быть найдено в классе функций,


которые могут иметь полюсы в тех точках плоскости $\zeta$, где они


имеются у функции $\omega(\cdot)$, а в случае круга --- еще и в бесконечно


удаленной точке. Полное исследование проводится таким же


образом, что и в (\cite{Mush.MTU}, \S\,125). В каждом конкретном случае


потенциалы будут зависеть от одной или от двух комплексных или


вещественных постоянных.


Ниже будет рассмотрен пример, а пока отметим следующее. В случае~1)


у функции $\omega(\cdot)$ в области $\Delta^+$ полюсов нет, если $D^+$ ---


ограниченная область. Если же область $D^+$ неограничена, то у функции


$\omega(\cdot)$ может быть простой полюс в $\Delta^+$ или на $\Lambda$.


В случае~2) бесконечно удаленная точка принадлежит линии $\Lambda$,


поэтому в области $\Delta^-$ у функции $\omega(\cdot)$ не может быть особых


точек.


Поэтому если область $D^+$ неограничена, то удобнее рассматривать в качестве


эквивалентной ей области полуплоскость $\Delta^-$.


В общем случае


$F(\cdot)=\mp F_0(\cdot)+G(\cdot)$, где $F_0(\cdot)$ --- интеграл типа Коши


с плотностью $f(\omega(\cdot))$, а $G(\cdot)$ --- некоторая рациональная


функция.





Существенно, что $\omega(\cdot)$ --- рациональная функция, т.\,к.


тогда и только тогда эта функция, определенная в области


$\Delta^{\pm}$, может быть естественным образом аналитически продолжена


через единичную окружность или через вещественную ось на всю комплексную


плоскость (за исключением, может быть, конечного числа точек, где у нее


расположены полюсы).





Рассмотрим пример. Пусть $D^-$ --- нижняя полуплоскость,


конформно эквивалентная единичному кругу $\Delta^+$ плоскости $\zeta$.


Отображающая функция имеет единственный простой полюс в точке $\zeta=1$


единичной окружности,


$$


\omega(\zeta)=i\frac{\zeta+1}{\zeta-1},\quad


\frac{\omega(\tau)}{\overline{\omega'(\tau)}}=


\frac{1}{2}\frac{\tau^2-1}{\tau^2}, \quad \tau\in\Lambda ,


$$


но у функции $F(\cdot)$ в этой точке нет особенности. Поэтому в решении


задачи о скачке (18) содержится только одна произвольная комплексная


постоянная. При $P=Q=0$


\begin{multline*}


F_0(\zeta) = -\frac{1}{2\pi i} \int_{ab} [\alpha(\tau)+\beta(\tau)]


\int_{\tau b}


\bigg[ \frac{1}{2\pi i}\int_{\Lambda}


\frac{1}{(\xi-\omega(x))(x-\zeta)} dx\bigg] d\xi\, d\tau + \\


%


+ \frac{1}{2\pi i} \int_{ab} [k\alpha(\tau)-\beta(\tau)]


\int_{\tau b}


\bigg[ \frac{1}{2\pi i}\int_{\Lambda}


\frac{1}{(\overline{\xi}-\overline{\omega(x)})(x-\zeta)} dx


\bigg] d\overline{\xi}\, d\tau - \\


%


- \frac{1}{2\pi i} \int_{ab} [\overline{\alpha(\tau)}


+\overline{\beta(\tau)}]


\bigg[ \frac{1}{2\pi i}\int_{\Lambda}


\frac{\tau-\omega(x)}{\overline{\tau}-\overline{\omega(x)}}


\frac{dx}{x-\zeta}


\bigg] d\overline{\tau} ,


\end{multline*}


все внутренние интегралы вычисляются по формуле Коши, если учесть, что


$\overline{\omega(1/\overline{\zeta})}=\omega(\zeta)$ и $1/\overline{x}=x$


при $x\in\Lambda$.


Дальнейшие вычисления проводятся аналогично случаю круга.


Как и следовало ожидать, вернувшись на плоскость $z$, получим выражения


вида (11) и (12) комплексных потенциалов для упругой полуплоскости


с дефектом.





С помощью потенциалов (17), как и в ранее рассмотренных случаях,


можно перейти от основных граничных задач для областей с дефектом,


конформно отображаемых на круг или полуплоскость, к сингулярным


интегральным уравнениям с логарифмической особенностью в ядре.





Как уже было отмечено, значения постоянных $P$, $Q$, $R$, $S$, $T$


в формулах (4), (5), (11), (12), (15), (16) и (17)


могут быть определены из каких-либо дополнительных условий.


Но, как следует из формул Колосова--Мусхелишвили, вычисленные по функциям


$\varphi(\cdot)$, $\psi(\cdot)$ напряжения и перемещения не зависят от этих


постоянных.





Численные решения интегральных уравнений с логарифмической особенностью


в ядре и сингулярных уравнений с ядром Коши в различных случаях


были получены методом Гал\"еркина (обсуждению результатов будет посвящена


другая работа). Если дефект расположен близко к границе или к особой


точке, то можно выделить ситуации, когда численные алгоритмы на основе


полученных в данной работе имеют преимущества перед другими.





\begin{thebibliography}{99}





\bibitem{Mush.MTU}


Мусхелишвили Н.И.


{\em Некоторые основные задачи математической теории упругости.


Основные уравнения. Плоская теория упругости. Кручение и изгиб}.


-- М.: Наука, 1966. -- 707~с.


\bibitem{PNB.Pre}


Плещинский Н.Б.


{\em Интегральные уравнения с логарифмической особенностью в ядре


для граничных задач теории упругости для плоскости, полуплоскости


и круга с дефектом вдоль гладкой дуги} //


Препринт \No\,1. Казанск. матем. об-во. -- Казань, 1997. -- 22~с.


\bibitem{Svr.DZUTT}


Саврук М.П.


{\em Двумерные задачи упругости для тел с трещинами}.


-- Киев: Наук. думка, 1981. -- 323~с.


\bibitem{PNB.PTIU}


Плещинский Н.Б.


{\em Приложения теории интегральных уравнений с логарифмическими


и степенными ядрами}. Учеб. пособие. --


Казань: Изд-во Казанск. ун-та, 1987. -- 157~с.


\bibitem{PNB.92}


Pleshchinskii N.B.


{\em Some classes of singular integral equations solvable


in a closed form and their applications} //


Pitman Research Notes in Math. Ser. Longman Scientific \&{}


Technical. -- New~York, 1991. -- V.256. -- P.246-256.


\bibitem{PP.IUTU}


Партон В.З., Перлин П.И.


{\em Интегральные уравнения теории упругост}и. --


М.: Наука, 1977. -- 312~с.


\bibitem{Gah.KZ.77}


Гахов Ф.Д. {\em Краевые задачи}. -- М.: Наука, 1977. -- 640~с.





\end{thebibliography}





\vskip 2cm





\post{\it Казанский государственный университет}{\it Поступила}


\post{}{12.11.1998}





\label{end}


\end{document}


